\section{Problem Statement}

We consider a distributed edge caching system in which a set of edge nodes collaboratively serve content requests generated by mobile users. Each node has limited cache capacity and may satisfy a request either locally, by retrieving the content from another node in the network, or by accessing the origin server. Since users move across nodes, the system exhibits mobility-induced correlations: nodes that frequently exchange users tend to observe related access patterns and may benefit from coordinated caching decisions.

We represent the system by a weighted graph \(G=(V,E,W)\), where \(V\) is the set of edge nodes, \(E\) is the set of connectivity relations, and \(W=\{w_{ij}\}\) captures mobility-induced interaction strength between nodes. Larger values of \(w_{ij}\) indicate stronger relationships, such as frequent handoffs or persistent overlap in user movement. Unlike classical caching systems in which each node is optimized independently, this setting gives rise to a coupled decision problem: content placement at one node affects accessibility, redundancy, and traffic at other nodes.

In practice, the network topology is not perfectly stable. Nodes may be added or removed due to mobility, intermittent availability, failures, or deployment changes. To model this uncertainty, we consider a perturbed graph \(\tilde{G}=(\tilde{V},\tilde{E},\tilde{W})\) sampled from a perturbation process \(\mathcal{P}(G)\), which captures random node arrivals and departures. A useful caching policy must therefore remain effective not only for the nominal graph \(G\), but also on average under such perturbations.

\subsection{Subproblems}

\subsubsection{Demand Modeling under Partial Observability}
The mobility trace provides user--node association information, but it does not include content-level request logs. As a result, demand must be synthesized while preserving realistic characteristics such as skewed popularity, spatial locality, and temporal variability. This creates a demand modeling problem: the request process should be realistic enough to support meaningful caching evaluation, yet controlled enough to avoid introducing artificial bias.

\subsubsection{Topology-Aware Cooperative Placement}
Caching decisions are coupled through the graph structure. If each node optimizes only for its own local demand, neighboring nodes may store the same content repeatedly, wasting cache space. On the other hand, fully centralized global optimization is costly and often unrealistic in distributed systems. The challenge is to exploit topology-aware cooperation without requiring complete global synchronization.

\subsubsection{Access--Replication Tradeoff}
Replicating content across more nodes improves accessibility and lowers retrieval latency, but it also consumes storage and incurs maintenance overhead. A good caching strategy must therefore balance the gain in access efficiency against the cost of storing and maintaining replicas.

\subsubsection{Relocation under Topology Perturbation}
When the topology changes due to node additions or removals, a previously efficient cache placement may become suboptimal. The system should be able to relocate cached content to adapt to perturbations, but relocation itself incurs communication and replacement cost. Thus, the policy must balance robustness against instability.

\subsection{Formal Problem Formulation}

Let \(\mathcal{C}\) denote the content catalog. For each node \(i \in \tilde{V}\) and content item \(c \in \mathcal{C}\), define the binary placement variable
\[
x_{i,c} =
\begin{cases}
1, & \text{if content } c \text{ is cached at node } i,\\
0, & \text{otherwise.}
\end{cases}
\]

Let \(x^{0}_{i,c}\) denote the initial cache state before perturbation. Let \(P_{i,c}\) denote the probability that node \(i\) requests content \(c\), and let \(f_c\) denote the replication-related cost of content \(c\), which may reflect update frequency, storage burden, or a combined maintenance factor.

Let \(d(i,j)\) be the shortest-path distance between nodes \(i\) and \(j\) in the perturbed graph \(\tilde{G}\), and let \(L(i,j)\) denote the latency of retrieving content from node \(j\) to node \(i\). We assume that \(L(i,j)\) is a non-decreasing function of graph distance. Let \(L_{\text{local}}\) denote the latency of serving content from the local cache, and let \(L_{\text{origin}}\) denote the latency of retrieving content from the origin server.

To simplify notation, define the best non-local retrieval latency for content \(c\) at node \(i\) under placement \(x\) and perturbed topology \(\tilde{G}\) as
\[
A_{i,c}(x,\tilde{G})
=
\min\!\left(
L_{\text{origin}},
\;
\min_{\substack{j \in \tilde{V} \\ j \neq i \\ x_{j,c}=1}} L(i,j)
\right).
\]
If no other node stores content \(c\), then \(A_{i,c}(x,\tilde{G}) = L_{\text{origin}}\).

The expected access cost is
\[
C_{\text{access}}(x,\tilde{G})
=
\sum_{i \in \tilde{V}} \sum_{c \in \mathcal{C}}
P_{i,c}
\left[
x_{i,c} L_{\text{local}}
+
(1-x_{i,c}) A_{i,c}(x,\tilde{G})
\right].
\]

The replication cost is
\[
C_{\text{rep}}(x,\tilde{G})
=
\sum_{i \in \tilde{V}} \sum_{c \in \mathcal{C}}
\gamma f_c x_{i,c},
\]
where \(\gamma > 0\) scales the importance of replication overhead.

To discourage redundant placement among strongly related nodes, we define the topology-aware redundancy penalty
\[
C_{\text{red}}(x,\tilde{G})
=
\sum_{c \in \mathcal{C}}
\sum_{\substack{i,j \in \tilde{V} \\ i<j}}
\tilde{w}_{ij}\,x_{i,c}x_{j,c}.
\]

To capture the cost of adapting the cache after perturbation, we define the relocation cost
\[
C_{\text{rel}}(x,x^{0})
=
\sum_{i \in \tilde{V}} \sum_{c \in \mathcal{C}}
\delta_c \, |x_{i,c} - x^{0}_{i,c}|,
\]
where \(\delta_c\) denotes the cost of adding or removing content \(c\) at a node.

The topology-aware cooperative caching problem is therefore formulated as the following stochastic optimization problem:
\[
\min_{x}
\;
\mathbb{E}_{\tilde{G} \sim \mathcal{P}(G)}
\left[
C_{\text{access}}(x,\tilde{G})
+ \lambda C_{\text{rep}}(x,\tilde{G})
+ \beta C_{\text{red}}(x,\tilde{G})
+ \mu C_{\text{rel}}(x,x^{0})
\right],
\]
where \(\lambda > 0\), \(\beta > 0\), and \(\mu > 0\) control the tradeoff between access efficiency, replication overhead, redundancy reduction, and relocation stability.

This objective captures the central design challenge of the system. Aggressive replication can reduce latency but increases replication and redundancy cost. Conservative placement may reduce overhead but increase expensive remote retrievals. Frequent relocation may adapt well to perturbations but can incur substantial update cost. The desired policy is therefore one that minimizes expected overall cost under random topology perturbations.

\subsection{Constraints}

The optimization problem is subject to the following practical constraints.

\paragraph{Cache Capacity Constraint.}
Each node has limited storage and therefore cannot cache the entire catalog. If \(K_i\) denotes the cache capacity of node \(i\), then the total size or number of cached items must remain within its budget:
\[
\sum_{c \in \mathcal{C}} x_{i,c} \le K_i,
\qquad \forall i \in \tilde{V}.
\]
This ensures that the solution is feasible under node-level storage limitations.

\paragraph{Binary Placement Constraint.}
Caching is a discrete decision: an item is either stored at a node or it is not. Therefore,
\[
x_{i,c} \in \{0,1\},
\qquad \forall i \in \tilde{V},\; \forall c \in \mathcal{C}.
\]
This makes the placement problem combinatorial.

\paragraph{Latency Ordering Constraint.}
The system assumes a natural service hierarchy: local cache access is cheapest, inter-node retrieval is more expensive, and origin retrieval is most expensive:
\[
L_{\text{local}} < L(i,j) \le L_{\text{origin}},
\qquad \forall i,j \in \tilde{V},\; i \neq j.
\]
This reflects the fact that cooperation is beneficial but not free.

\paragraph{Perturbation Feasibility Constraint.}
Cache placement after perturbation is only defined on nodes that remain present in \(\tilde{V}\). Newly added nodes may begin with empty cache state, while removed nodes cannot participate in storage or retrieval. This ensures that placement and relocation decisions are consistent with the realized perturbed topology.

\paragraph{Topology-Dependent Communication Constraint.}
Inter-node retrieval cost depends on the perturbed graph structure through \(L(i,j)\). Thus, cooperation is not assumed to be globally uniform or free; it is mediated by graph connectivity and path cost. This ensures that the topology plays a direct and meaningful role in the optimization problem.